\label{chap:background}

The \gls{CL} is a procedure that allows doctors to examine the inner lining of the colon and rectum using a flexible tube with a lighted camera. The \gls{CLP} passes through the anus and rectum into the colon, providing images of the interior. Those images can be used to identify abnormalities such as polyps, which are growths that can develop into colorectal cancer if left untreated. 

However, the accuracy of colonoscopy can be affected by various factors, as discussed in Introduction \ref{chap:introduction}. Considering the importance of early detection of colorectal cancer, is undeniable that improving the accuracy of colonoscopy is a critical area of research. 

\section{Machine Learning and Computer Vision}
\label{sec_machine_learning}

In 1943 the first mathematical model of what would be the bases to \gls{AI},field of computer science dedicated to creating systems capable of performing tasks that typically require human intelligence, was published by Warren McCulloch and Walter Pitts in \cite{mcculloch1943} paper. The article evalueats the neuron activity as a character of "all or none", where their events can be treated by propositional logic. Biologically, the brain cell has needs to receive an excitation higher than it`s thredhold to trigger an impulse, which is determined by their own characteristcs.

Considering that, the authors claim that an neuron can be described as a logic block, executing logical function. Brain cells could be organized in nets that could describe any logical preposition. With the question 'Can machines think?', Alan Turing starts his paper \cite{DBLP:journals/x/Turing50}, describing thinking such as examples of such tasks includes problem-solving, decision-making, and pattern recognition. In this approach, intelligence could be simulated by computational process.


\gls{ML} is an area of knowledge that focus on developing algorithms that allow computer system to learn from data and improve their performance without being explicit programmed. This mathematical relationship derived from the data is called model. Each revelant feature, or characteristc of data, must be represented numerically, so the model can learn. In \cite{5392560}, Samuel, A. L shows a computer that learns to play a game called checker better than the person who wrote the program. It was chosen because of the simplicity of its rules and transparency of the computer learning progress. The machine should look ahead a few moves and evalueats the result board position. The score of the result position was in terms of linear polynomial with a total of 38 parameters, but evalueats only 16 at a time for a question of efficiency and relevance. The algorithm should find out which of those parameters was relevant. The conclusion of the 1959`s article was possible that devise learning schemes would outperform an average person and those schemes should be applied to real-life problems.

\begin{figure}[htbp]
    \centering
    \includegraphics[width=0.8\textwidth]{checker}
    \caption{Eight move opening in checker by machine learning model. This image is derivated from Samuel paper \cite{5392560}.
    \label{Eight move opening in checker by machine learning model.}}
\end{figure}



\section{Data organization and Datasets}
\label{sec_dataset}
Not only in \gls{AI} context, data organization is more important than ever in the XXI century, considering that companies collect large amounts of data from various sources daily. Considering that, data must often be organized collections that make information accessible for analysis and application - called \gls{DS}.
However, there are more than one way of storing data, and different types of \gls{DS} store it in various ways. 

Structured \gls{DS} organize information in predefined structure, normally tables which each column is an attribute and each row an instance. In other hand, unstructured \gls{DS} contain information that does not fit in a rigid schema or cannot be resumed to a numerical or text value. Some examples of it can be images, videos, web-pages, documents or even applications. Futhermore, semistructured \gls{DS} don't follow rigid schemas, they incorporate defined syntax, such as \gls{JSON}, \gls{HTML} or \gls{XML}.

The type of task that can be solved or explored using \gls{ML} and \gls{AI} depends not only in the quality of data, but the type of the \gls{DS} that will train the model.

\section{Supervised Learning}
\label{sec_supervised_learning}
Frank Rosenblatt work \cite{rosenblatt1958perceptron} proposed in 1958 a model to explain how the human brain storages data and utilize it for pattern recognition and behavior. The probabilistic model called perceptron is divided in three layers: Sensority units that respondes to impulses in all or nothing bases; Association units receives the transmitted impulses and, if the algebric sum of those are greated than the unit thredhold, it fires and generates the responses of some Response units. Each response unit has feedbacks connections that can be excitatory or inhibitory and are mutually exclusive.If a response occurs, it will inhibit others.

\begin{figure}[htbp]
    \centering
    \includegraphics[width=0.8\textwidth]{perceptron}
    \caption{Probabilistc model for information storage organization displayed in \cite{rosenblatt1958perceptron}.
}
\end{figure}

A stimulus can be classified in classes, which provokes a stronger impulse of it corresponding response. However, those responses are not based on boolean algebra, but probabilistic logic. The system is capable of modifying association units or their connections based on the stimulus and impulses provided. Frank calls "differentiated environment" where each response is associeated to a distinct class of similar stimuli, where performance can be enhanced by the use of binary system.

Bernard Widrow \cite{widrow1960adaptive} develop a system that learns by giving a lot of examples of inputs and their respective desired outputs, instead of comparing every logic detail. The quality of the system performance is directly related to the number of examples given. To illustrate it, a machine called \gls{ADALINE} was built. Similar to von Neumann \cite{Neumann1956ProbabilisticLA} with his neuron built as a combinatorial logic circuit, each input is linearly combined with it respectives weights, which are setted during the training procedure.
\begin{figure}[htbp]
    \centering
    \includegraphics[width=0.6\textwidth]{vonNeumein-neuron}
    \caption{Adjustable neuron cited by Bernard Window \cite{widrow1960adaptive}.
    }
\end{figure}
\begin{figure}[htbp]
    \centering
    \includegraphics[width=0.6\textwidth]{adaline-schematic}
    \caption{Adaline schematic cited by Bernard Window \cite{widrow1960adaptive}.
    }
\end{figure}
\gls{ADALINE} learning is stored at the weights a0...a16 and can be used to classify the original patterns and noisy it`s distorsed versions. The training method is an iterative search process to find the best gain set, which can be look by the error of an instance
\begin{equation}
    \epsilon(k) = d(k) - [a_0 + \sum_{j=1}^{n} a_j s_j(k)]
\end{equation}

Widrow and Holf demonstrated mathematically that the use \gls{MSE} to adaptation weights was because of it parabolic propertys of single global minimum.
\begin{equation} 
    MSE = \frac{1}{m} \sum_{i=1}^{m} (\hat{y}^{(i)} - y^{(i)})^2 
\end{equation}
With that, minimizing \gls{MSE} is equivalent to minimizing the average number of neuron error, which lead to the best set of weights. The training process uses partial derivatives of the formula above, showing that the direction and the magnitude of each weight depends exclusively from the input signal and the magnitude error, as detailed in \ref{sub_sec_backprop_gradi}.
\begin{equation} 
\frac{\partial \epsilon^2(k)}{\partial a_j} = -2s_j(k)\epsilon(k) 
\end{equation}

\section{Artificial Neural Networks and Deep Learning}
\label{sec_deep_learning}
For classification tasks is very common to use \gls{ANN} algorithms, a \gls{ML} model based on the biological neurons for \gls{SL}. In the human brain, neurons are connect to each other and communicate through electrical signals. Mathematically, this structure is represented by a fuction \(f(X)\) , that maps an input \(X = x1, x2, x3, ...\) to predict a output \(Y\).

\gls{ANN} is composed of three different parts, called \textit{layers}. Each one being composed of a group of \textit{neurons}. The first one is the \textit{input layer}, which receives raw data. The second consist of multiple layers, the \textit{hidden layer} consist of \textit{nodes} that transform inputs into new representation. They can be mathematically represented as a \textbf{linear transformations}, which the input features are multiplied by their associated weights and added bias to pass from one layer to the another. Finally, the \textit{output layer} applys to the output of the hidden layer a \textbf{nonlinear activation function} such as sigmoid, Relu, tanh, softmax and produce the final prediction.

Models with at least 4 layers are called \gls{DL} models. Not only can data be splitted into smaller pieces to make predictions, \gls{DL} often needs a massive amount of training data and computer perfomrance to reach optimal performance, considering the fact that some of more than millions or even a billion of parameters. The inclusion of multiple \textit{hidden layers} distinguishes a \gls{DL} from a "simple" \gls{ANN}, which allows the model to learn complex hierarchical features of data. The first ones identify broader paterns such as borders and deeper ones identify more granular patterns.

\subsection{Backpropagation and Gradient Descendent}
\label{sub_sec_backprop_gradi}
A \gls{DL} model updates its parameters in order to achieve greater precision in output predictions. To do so, is necessary that each of one of the millions, or even billions, parameters to be calculed separately, what can be have and tremendous computational cost considering the amount of neurons in a network. This process occurs after each \textbf{epoch} - a complete pass throught the train dataset, where the model process every sample, calculates loss and, finally, update it`s parameters.

Backward propagation erro, also referenced as \gls{BP}, is an algorithm to calculate how changes an individual weight or bias in  \gls{ANN} will affect the accuracy of predictions. By calculating the \textbf{gradient} - vector of partial derivatives of the \gls{LF} for each variable in every equation, it shows optimization algorithms - technics to adjust model parameters, each equations and parameters should be adjusted to reduce \gls{LF}.


\gls{GD} is an optimization algorithm  for minimizing a differentiable multivariate function. In \gls{SL}, it commonly works along side \gls{BP}, based on \textit{learning rate} - size of the steps necessary to reach minimum gradient. High ones means in highers steps, but may skip de converge point. However, small ones kill optimization efficiency.

For big amounts of data, dataset may be splitted into \textit{batchs} for efficiency. In this approach, an \textbf{epoch} consist of multiple iterations, which each process one batch. The number of iterations can be calculated by: \( iterationPerEpoch = totalSamples / batchSize\).

\subsection{Stochastic Gradient Descendent}
\label{sub_sec_stoch_grad}
Based on \gls{GD}, \gls{SGD} executes and updates parameters for each sample in dataset, one at a time inside an epoch. Easier to stay in memory than a \textit{batch} approach, it can reduce computational efficiency and generate noisy gradients.

\subsection{Mini-Batch Gradient Descendent}
\label{sub_sec_mini_batch_grad}
It combines \gls{SGD} with \textit{batch} approach, splitting dataset in small batches, called \textit{mini-batches} and updates model parameters in each of them.

\subsection{SGD with Momentum}
\label{sub_sec_sgd_momentum}
\gls{GDM} is an extension of \gls{GD} algorithm. Based on physics, where movement depends on current force and past velocity. The main goal of this approach is to smooth learning processo by adding past gradients, what leads to more stable convergence, which avoids \gls{BP} to skip the global minimum.

A new parameter called \textit{momentum factor} controls how mmuch of the past gradients are "remembered" in a update. Zero (0) means the optimizer will not consider past gradients and a value closes to one (1) will consider more past values.


\subsection{Adaptive Moment Estimation}
\label{sub_sec_adam}
Adaptive Moment Estimation, or \gls{ADM}, is an optimization algorithm such as \gls{GD}. First appear in \cite{DBLP:journals/corr/KingmaB14} Diederik P. Kingma and Jimmy Ba show that 


\subsection{Adamax}
\label{sub_sec_adamax}


\section{Models}
\label{sec_models}

\section{Evaluated Metrics}
\label{sec_metrics}